\section{Formalizacja agenta bazowego} \label{sec:formalization}

Sekcja zawiera formalny opis agenta parametrycznego, będącego podstawą wszystkich modyfikacji. Notację ujednolicono z opisem rozszerzeń przedstawionych w dalszej części pracy. Notacja różni się od notacji stosowanej przez García-Sánchez i in.~\cite{sanchez2020}.

\subsection{Reprezentacja i strategia decyzyjna}

W poniższych wzorach indeks górny $(\cdot)^{\text{wł}}$ odnosi się do gracza, a $(\cdot)^{\text{prz}}$ --- do przeciwnika. Agent przechowuje wektor wag $\mathbf{w} \in [0,1]^D$, gdzie $D$ jest liczbą parametrów danego wariantu. W każdej turze, dla każdej legalnej akcji $a$ dostępnej w stanie $s$, agent symuluje stan wynikowy $s' = T(s, a)$ i oblicza wartość akcji:

\begin{equation}
    f(s, a, \mathbf{w}) =
        \Delta H^{\text{prz}} - \Delta H^{\text{wł}}
        + \Delta M^{\text{prz}} - \Delta M^{\text{wł}}
        + \Delta S^{\text{prz}} - \Delta S^{\text{wł}}
        - \Delta \text{mana}
    \label{eq:score}
\end{equation}

Następnie agent wybiera akcję o najwyższej wartości:
\begin{equation}
    a^* = \arg\max_{a \in \mathcal{A}(s)}\; f(s, a, \mathbf{w})
    \label{eq:greedy}
\end{equation}
gdzie $\mathcal{A}(s)$ oznacza zbiór wszystkich legalnych akcji dostępnych
w stanie $s$ --- jest to strategia zachłanna z horyzontem 1. Wyjątkiem od równania \eqref{eq:greedy} są stany terminalne: akcja prowadząca do wygranej (punkty zdrowia bohatera przeciwnika $\leq 0$) otrzymuje wartość $+\infty$, a prowadząca do przegranej (własne punkty zdrowia $\leq 0$) --- wartość $-\infty$.

\subsection{Składniki funkcji oceny}

Każdy składnik równania~\eqref{eq:score} opisuje \textbf{różnicę} pewnej cechy między stanem przed akcją ($s$) a po akcji ($s'$), zawsze z perspektywy gracza, którego cecha ta dotyczy. Wartości dodatnie są korzystne z punktu widzenia gracza wykonującego akcję (tj. wzrost $\Delta H^{\text{prz}}$ oznacza, że przeciwnik utracił punkty zdrowia, co jest pożądane).

\paragraph{Zmiana stanu bohatera $\Delta H$.}

Dla bohatera gracza $p$ przed akcją i $p'$ po akcji:
\begin{equation}
    \Delta H(p, p', \mathbf{w}) =
        w_{\text{HP}} \cdot
        \bigl[(z_p + \text{pan}_p) - (z_{p'} + \text{pan}_{p'})\bigr]
        + w_{\text{HA}} \cdot
        \bigl[\text{atk}_p - \text{atk}_{p'}\bigr]
    \label{eq:delta_h}
\end{equation}
gdzie $z_p$ oznacza punkty zdrowia bohatera, $\text{pan}_p$ --- jego pancerz (ang.~\textit{armor}), a $\text{atk}_p$ --- aktualny atak bohatera. Dodatni wynik oznacza, że bohater utracił punkty, co z perspektywy \emph{przeciwnika} jest korzystne.

\paragraph{Funkcja wartości stronnika $\nu$.}

Przed obliczeniem zmiany stanu planszy definiuje się pomocniczą funkcję wartości pojedynczego stronnika $m$. Wprowadza się następujące oznaczenia:
\begin{itemize}
    \item $\text{zdr}(m)$ --- punkty zdrowia stronnika,
    \item $\text{atk}(m)$ --- atak stronnika,
    \item $\text{koszt}(m)$ --- koszt many karty stronnika,
    \item $r(m) \in \{0, 1, 2, 3, 4\}$ --- rzadkość karty (brak/zwykła, pospolita, rzadka, epicka, legendarna),
    \item $\mathbf{1}[\text{posiada}(m, k)]$ --- indykator wskazujący, czy stronnik $m$ posiada zdolność $k$.
\end{itemize}

\begin{equation}
    \nu(m, \mathbf{w}) =
        w_{\text{mH}} \cdot \text{zdr}(m)
        + w_{\text{mA}} \cdot \text{atk}(m)
        + \sum_{k \in \mathcal{K}} w_k \cdot \mathbf{1}[\text{posiada}(m, k)]
        + w_{\text{cost}} \cdot \text{koszt}(m)
        + w_{\text{rar}} \cdot r(m)
    \label{eq:minion_value}
\end{equation}

Zbiór $\mathcal{K}$ zawiera dziewięć zdolności specjalnych stronnika:
\[
\mathcal{K} =
\left\{
\begin{aligned}
&\text{Szarża},\; \text{Agonia},\; \text{Boska tarcza},\; \text{Inspiracja},\; \text{Kradzież Życia},\\
&\text{Ukrycie},\; \text{Prowokacja},\; \text{Furia Wichru},\; \text{Trucizna}
\end{aligned}
\right\}
\]

\paragraph{Zmiana stanu planszy $\Delta M$.}

Niech $B$ i $B'$ oznaczają strefę stronników gracza przed i po akcji. Wyróżnia się trzy klasy zdarzeń:
\begin{itemize}
    \item stronnik $m \in B \cap B'$ przeżył --- zmiana zdrowia lub ataku,
    \item stronnik $m \in B \setminus B'$ zginął,
    \item stronnik $m \in B' \setminus B$ pojawił się na planszy.
\end{itemize}

\begin{equation}
\begin{split}
    \Delta M(B, B', \mathbf{w}) =\;
        & \sum_{m \in B \cap B'} \Bigl[
            w_{\text{HR}} \cdot \Delta\text{zdr}(m) \cdot \nu(m, \mathbf{w})
            + w_{\text{AR}} \cdot \Delta\text{atk}(m) \cdot \nu(m, \mathbf{w})
          \Bigr] \\
        +\; & \sum_{m \in B \setminus B'} w_{\text{kill}} \cdot \nu(m, \mathbf{w}) \\
        -\; & \sum_{m \in B' \setminus B} w_{\text{app}} \cdot \nu(m, \mathbf{w})
\end{split}
    \label{eq:delta_m}
\end{equation}

gdzie $\Delta\text{zdr}(m) = \text{zdr}_B(m) - \text{zdr}_{B'}(m)$ oraz analogicznie $\Delta\text{atk}(m) = \text{atk}_B(m) - \text{atk}_{B'}(m)$. Trzeci człon jest odejmowany, ponieważ pojawienie się nowego stronnika przeciwnika jest zdarzeniem niekorzystnym.

\paragraph{Zmiana sekretów $\Delta S$.}

\begin{equation}
    \Delta S(p, p', \mathbf{w}) =
        w_S \cdot \bigl(|Z^{\text{sek}}_p| - |Z^{\text{sek}}_{p'}|\bigr)
    \label{eq:delta_s}
\end{equation}
gdzie $|Z^{\text{sek}}_p|$ oznacza liczbę aktywnych sekretów gracza $p$.

\paragraph{Koszt many $\Delta\text{mana}$.}

\begin{equation}
    \Delta\text{mana}(s, s', \mathbf{w}) =
        w_{\text{mana}} \cdot
        \bigl(\text{mana}^{\text{rem}}_s - \text{mana}^{\text{rem}}_{s'}\bigr)
    \label{eq:delta_mana}
\end{equation}
gdzie $\text{mana}^{\text{rem}}$ oznacza pozostałą ilość many gracza. Składnik ten jest \emph{odejmowany} w równaniu~\eqref{eq:score}, penalizując akcje o wysokim koszcie many.

\subsection{Wektor wag i genom}

Bazowy agent posiada $D = 21$ parametrów. Pełny wektor wag $\mathbf{w} \in [0,1]^{21}$ zestawiono w tabeli~\ref{tab:weights21}. Wektor $\mathbf{w}$ stanowi \textbf{genom} osobnika w algorytmie ewolucyjnym. Przestrzeń przeszukiwań to $[0,1]^D$, która jest wyznaczana przez każdego agenta.

\begin{table}[h]
\centering
\caption{Parametry agenta bazowego (21 wag).}
\label{tab:weights21}
\begin{tabular}{cll}
\toprule
Indeks & Oznaczenie & Opis \\
\midrule
1  & $w_{\text{HP}}$            & Redukcja zdrowia bohatera \\
2  & $w_{\text{HA}}$            & Redukcja ataku bohatera \\
3  & $w_{\text{HR}}$            & Redukcja zdrowia stronnika \\
4  & $w_{\text{AR}}$            & Redukcja ataku stronnika \\
5  & $w_{\text{app}}$           & Pojawienie się stronnika na planszy \\
6  & $w_{\text{kill}}$          & Zabicie stronnika \\
7  & $w_S$                      & Usunięcie sekretu \\
8  & $w_{\text{mana}}$          & Zużyta mana \\
9  & $w_{\text{mH}}$            & Zdrowie stronnika (wartość ogólna) \\
10  & $w_{\text{mA}}$           & Atak stronnika (wartość ogólna) \\
11 & $w_{\text{Charge}}$        & Zdolność: Szarża \\
12 & $w_{\text{Deathrattle}}$   & Zdolność: Agonia \\
13 & $w_{\text{DivineShield}}$  & Zdolność: Boska Tarcza \\
14 & $w_{\text{Inspire}}$       & Zdolność: Inspiracja \\
15 & $w_{\text{LifeSteal}}$     & Zdolność: Kradzież Życia \\
16 & $w_{\text{Stealth}}$       & Zdolność: Ukrycie \\
17 & $w_{\text{Taunt}}$         & Zdolność: Prowokacja \\
18 & $w_{\text{Windfury}}$      & Zdolność: Furia Wichru \\
19 & $w_{\text{Poisonous}}$     & Zdolność: Trucizna \\
20 & $w_{\text{rar}}$           & Rzadkość stronnika \\
21 & $w_{\text{cost}}$          & Koszt stronnika \\

\bottomrule
\end{tabular}
\end{table}
